\documentclass[11pt]{book}
\usepackage[margin=1in]{geometry}
\usepackage[T1]{fontenc}
\usepackage{hyperref}
\usepackage{longtable}
\usepackage{listings}
\usepackage{xcolor}
\lstdefinestyle{code}{
  basicstyle=\ttfamily\small,
  breaklines=true,
  frame=single,
  columns=fullflexible
}
\title{DuckDB Predictive Agent\\Implementation and Extension Guide}
\author{Codex MVP Scaffold}
\date{\today}

\begin{document}
\maketitle
\tableofcontents

\chapter{Purpose}
This document explains the implementation scaffold for the DuckDB Predictive Agent MVP. It is
written as a practical build book: what exists, why it is structured this way, how to extend it,
and where the current boundaries are.

\chapter{System Overview}
The product is a safe predictive reasoning system over uploaded DuckDB databases. The core safety
contract is:

\begin{quote}
LLM proposes. Backend validates. Backend executes. Backend summarizes. LLM reasons.
\end{quote}

The scaffold currently includes a heuristic fallback path so the product remains runnable even
before live LLM integration is enabled.

\section{Top-Level Modules}
\begin{longtable}{|p{0.24\textwidth}|p{0.68\textwidth}|}
\hline
Module & Responsibility \\
\hline
\texttt{backend/app/api} & FastAPI routes for uploads, connections, profiles, predictions, and logs \\
\hline
\texttt{backend/app/db} & SQLAlchemy metadata models and session lifecycle \\
\hline
\texttt{backend/app/duckdb\_tools} & Read-only connector, profiler, and SQL safety layer \\
\hline
\texttt{backend/app/memory} & Database-scoped semantic memory persistence and retrieval \\
\hline
\texttt{backend/app/agent} & Planning, exploration, evidence, prompts, and prediction boundaries \\
\hline
\texttt{backend/app/llm} & Provider configuration and hosted/local backend switching layer \\
\hline
\texttt{backend/app/services} & Application orchestration for uploads, databases, and predictions \\
\hline
\texttt{frontend} & Streamlit demo UI \\
\hline
\end{longtable}

\chapter{Backend Architecture}
\section{Configuration}
Runtime settings live in \texttt{backend/app/config.py}. Model selection lives in a separate LLM
config file so business logic does not hard-code model names. This is the main extension seam for:

\begin{itemize}
\item switching from OpenAI to a local runtime
\item assigning different models to planning, exploration, and final prediction
\item adding fallback and summary-specific models
\end{itemize}

\section{Metadata Store}
SQLite is used for MVP speed and simplicity, but the application layer is written with SQLAlchemy
so a later Postgres migration is straightforward. The core metadata entities are:

\begin{itemize}
\item users
\item database\_connections
\item table\_profiles
\item column\_profiles
\item relationship\_profiles
\item semantic\_memories
\item prediction\_runs
\item query\_audit\_logs
\end{itemize}

\section{Upload and Connection Flow}
The MVP upload flow accepts a DuckDB file, stores it under managed app storage, and registers a
database connection record. This avoids unsafe arbitrary path access. Profiling runs immediately
after connection creation.

\section{Profiling}
The profiler performs deterministic inspection:

\begin{itemize}
\item table discovery
\item column discovery
\item row counts
\item null and distinct counts
\item numeric, temporal, and categorical summaries
\item primary key heuristics
\item relationship heuristics
\item sample rows
\end{itemize}

This deterministic profile is the base truth layer. LLM summaries should only consume this
controlled representation, not unrestricted raw database contents.

\section{SQL Safety}
The SQL safety layer is deliberately isolated in \texttt{sql\_safety.py}. It combines:

\begin{itemize}
\item parser-based validation through \texttt{sqlglot}
\item a secondary forbidden-token gate
\item automatic row limiting
\item process-based execution timeout boundaries
\end{itemize}

This module is one of the most important future hardening surfaces. A production system should
extend it with stronger schema-aware validation and tighter worker isolation.

\chapter{Prediction Pipeline}
\section{Current Runtime Path}
The implementation includes both the intended modular LLM boundaries and a runnable heuristic
fallback. Today, the prediction service:

\begin{enumerate}
\item loads the profiled database summary and memory
\item derives a heuristic predictive task when live LLMs are disabled
\item generates safe exploratory queries
\item validates and executes them through the safety layer
\item records query audits
\item builds an evidence packet
\item returns a heuristic prediction summary
\end{enumerate}

\section{LLM Boundaries}
The future LLM-backed path should activate these components:

\begin{itemize}
\item \texttt{TaskPlanner}
\item \texttt{SQLExplorer}
\item \texttt{FinalPredictor}
\end{itemize}

Each component expects structured JSON outputs and should remain independently testable.

\section{Evidence Packet}
The evidence accumulator centralizes the packet that will ultimately be shown to users and passed
into the final predictor. This object is intentionally stable because many future capabilities
depend on it:

\begin{itemize}
\item confidence scoring
\item explanation generation
\item UI evidence rendering
\item offline evaluation
\item tracing and debugging
\end{itemize}

\chapter{Frontend}
The Streamlit frontend is intentionally simple. Its job is not product polish; it is to make the
workflow observable while backend behavior is being validated. It supports:

\begin{itemize}
\item file upload
\item database connection
\item profile display
\item predictive question submission
\item answer display
\item query-log display
\end{itemize}

\chapter{Extension Opportunities}
\section{Hosted and Local LLMs}
The next major extension is to implement the local provider backend. The configuration layer already
anticipates a future \texttt{model\_path} field for local models such as a Llama 70B deployment.
Reasonable backend options include:

\begin{itemize}
\item \texttt{llama.cpp}
\item \texttt{vLLM}
\item an internal REST wrapper around a local inference server
\end{itemize}

\section{Confidence Engine}
Today confidence is heuristic. A stronger implementation should compute confidence from evidence
quality signals such as:

\begin{itemize}
\item table coverage
\item key target-signal availability
\item temporal coverage quality
\item successful query ratio
\item entity-level support strength
\end{itemize}

\section{Query Planning}
The current heuristic query generator is intentionally conservative. It should later be replaced or
augmented with:

\begin{itemize}
\item planner-guided query selection
\item schema-aware join planning
\item result summarization compression
\item reusable prediction recipes
\end{itemize}

\section{Critic-Guided Iterative Evidence Acquisition}
One important extension is to add a critic that inspects the final prediction prompt, the final
predictor output, and the ground truth for benchmarked tasks. The critic should answer a narrower
question than ``was the final answer wrong?'': it should decide whether the final prompt contained
enough information for the predictor to recover the ground truth.

This design is useful because it creates a cleaner signal for improving the planner. If the prompt
was insufficient, the planner and exploration stages likely failed to retrieve enough evidence. If
the prompt was sufficient but the answer was still wrong, the final predictor is the more likely
failure point.

\subsection*{Core Loop}
\begin{enumerate}
\item the solver or task planner creates an initial predictive plan
\item the explorer gathers evidence from the database through safe SQL
\item the prompt builder composes the final evidence prompt
\item the final predictor produces an answer
\item the critic checks prompt sufficiency against ground truth
\item if the prompt is insufficient, the critic requests additional evidence types
\item the solver or explorer fetches follow-up evidence
\item the prompt builder integrates the new evidence and rebuilds the prompt
\item the final predictor retries
\end{enumerate}

\subsection*{Why The Critic Should Request Evidence Types}
The critic should request missing evidence categories rather than raw SQL. For example, it may ask
for recent support complaints, failed payments in a time window, or inactivity-based recency
signals. Query generation should still remain the responsibility of the solver or explorer so the
safety boundary stays intact.

\subsection*{Prompt Builder Requirements}
The prompt builder should not assume a fixed evidence shape. It should support:

\begin{itemize}
\item modular evidence sections
\item dynamic insertion of critic-requested evidence
\item deduplication of repeated facts
\item prioritization of high-value evidence under context limits
\item provenance tags indicating whether evidence was part of the initial plan or a critic-guided refinement
\end{itemize}

A practical evidence layout can include:

\begin{itemize}
\item core facts
\item candidate entities
\item temporal signals
\item payment signals
\item support or friction signals
\item critic-requested evidence
\item limitations
\item queries run
\end{itemize}

\subsection*{Training Signal For Planner Improvement}
This loop also creates strong training data for the planner. Each trajectory can store:

\begin{itemize}
\item the original planner output
\item the critic sufficiency judgment
\item the critic-requested additional evidence types
\item the refined planner or exploration behavior
\item the final benchmark score against ground truth
\end{itemize}

These traces can support supervised fine-tuning, preference optimization, or reward-based training
for the planner without relying on fully manual labels.

\section{Adapter-Based Fine-Tuning}
The preferred fine-tuning mechanism for this project is no longer a single monolithic model update.
Instead, the training system assumes:

\begin{itemize}
\item one shared base model
\item one adapter for \texttt{task\_planner}
\item one adapter for \texttt{sql\_explorer}
\item one adapter for \texttt{final\_predictor}
\item one adapter for \texttt{critic}
\end{itemize}

This gives the project a practical compromise between serving cost and specialization. The base
model remains shared across roles, while the role-specific behavior lives in adapters that can be
trained independently.

\subsection*{Why Adapters Fit This System}
Adapters are a strong fit here because the four main agent roles share the same broad reasoning
capability but differ in objective:

\begin{itemize}
\item the planner structures the task
\item the explorer chooses safe evidence-gathering queries
\item the predictor synthesizes the final answer
\item the critic judges sufficiency and requests follow-up evidence
\end{itemize}

Using adapters allows one role to improve without causing broad regressions in the others.

\subsection*{Role-Sampled Training Episodes}
The fine-tuning loop should sample one training role per episode. A typical episode now becomes:

\begin{enumerate}
\item randomly choose a dataset
\item randomly choose a task from that dataset
\item randomly choose an entity from the train split
\item randomly choose a trainable role
\item activate that role's adapter conceptually for the episode
\item run the full trajectory with the normal solver, explorer, predictor, and critic stages
\item assign responsibility using the critic and benchmark score
\item export a role-specific training example for the selected adapter
\end{enumerate}

The shared base model and all non-selected adapters are treated as fixed for that training episode.

\subsection*{Role-Specific Training Artifacts}
The implementation should export separate JSONL files for each role, for example:

\begin{itemize}
\item \texttt{task\_planner\_training\_examples.jsonl}
\item \texttt{sql\_explorer\_training\_examples.jsonl}
\item \texttt{final\_predictor\_training\_examples.jsonl}
\item \texttt{critic\_training\_examples.jsonl}
\end{itemize}

These files are intended for later cluster-side training jobs and should include metadata about:

\begin{itemize}
\item target role
\item target adapter
\item shared base model
\item dataset and task identity
\item entity identity
\item benchmark score
\end{itemize}

\subsection*{Operational Consequence}
At inference time, the ideal production setup is:

\begin{itemize}
\item load one shared base model
\item activate the planner adapter for planning
\item switch to the explorer adapter for SQL generation
\item switch to the predictor adapter for final answering
\item switch to the critic adapter for critique
\end{itemize}

This preserves modular behavior without paying the full cost of loading multiple unrelated models.

\subsection*{Local Runtime Resolution}
The local runtime now supports a concrete resolution rule for parameter-size aliases such as
\texttt{1b}, \texttt{3b}, \texttt{8b}, and later \texttt{70b}. For a local role call:

\begin{enumerate}
\item resolve the role to a size alias such as \texttt{8b}
\item map that alias to the configured base model path
\item check whether a role adapter exists under \texttt{finetuned\_llm/<size>/<role>}
\item if present, attach that adapter to the already loaded base model
\item if \texttt{finetuned\_llm/<size>/config.json} exists, allow that directory to act as a merged fine-tuned model copy
\end{enumerate}

This gives the runtime two compatible deployment styles:

\begin{itemize}
\item base model plus role adapters
\item merged fine-tuned model copy
\end{itemize}

\subsection*{Single In-Memory Base Load}
The most important runtime property is that the local system should not load the same base model
multiple times when multiple roles share the same underlying model path. Instead:

\begin{itemize}
\item cache one loaded base model per resolved base path
\item reuse that in-memory base model across planner, explorer, predictor, and critic steps
\item attach or switch role adapters on top of that shared base instance
\end{itemize}

This keeps the architecture aligned with the intended operational goal:

\begin{quote}
one shared local base model in memory, multiple role-specific adapters layered on top
\end{quote}

The practical consequence is that if all four roles resolve to the same \texttt{8b} base path, the
runtime should load that base model once and then reuse it throughout the full agent pipeline.

\section{Profiler Enrichment}
Useful profiler upgrades include:

\begin{itemize}
\item better foreign-key inference
\item column semantic typing
\item table role classification
\item anomaly and sparsity detection
\item incremental profile refresh support
\end{itemize}

\chapter{Operational Guidance}
\section{Running the Stack}
\begin{lstlisting}[style=code,language=bash]
uvicorn backend.app.main:app --reload
streamlit run frontend/streamlit_app.py
\end{lstlisting}

\section{Testing}
\begin{lstlisting}[style=code,language=bash]
pytest
\end{lstlisting}

\section{Configuration}
Copy the example config file and adapt it:

\begin{lstlisting}[style=code]
config/llm_config.example.json -> config/llm_config.json
\end{lstlisting}

\chapter{Known Gaps}
The scaffold is intentionally substantial but not feature-complete. The main remaining work items
are:

\begin{itemize}
\item live LLM execution path wiring into the prediction service
\item richer ranked entity generation
\item deletion endpoints and cascade handling
\item full Alembic migration setup
\item stronger sandboxing for production deployment
\item richer sample datasets and integration fixtures
\item critic-guided online refinement and planner improvement loops
\item cluster-side adapter training and adapter switching at inference time
\end{itemize}

\chapter{Conclusion}
This codebase is now structured to support incremental implementation without rework. The modules
separate concerns cleanly enough that profiling, safety, memory, and model orchestration can evolve
independently while preserving the core product contract.

\end{document}
